\section{Implementação e modelo temporal}
\label{sec:implementacao}

\subsection{Arquitetura e operações simuladas}

A aplicação foi organizada em quatro tarefas funcionais e duas auxiliares. A Figura~\ref{fig:arquitetura} apresenta as filas e o estado compartilhado. A implementação utiliza os mecanismos de tarefas, filas e sincronização do FreeRTOS; os cálculos de fusão, PID e mistura de motores estão escritos no código da aplicação. O Python foi utilizado somente após as coletas, para organizar dados e produzir figuras.

\figura{12_arquitetura.pdf}{Arquitetura do autopiloto didático.}{fig:arquitetura}{A medição de segurança termina após a carga artificial; a medição da telemetria termina antes da transmissão do texto.}

A \codigo{FUS_IMU} executa nominalmente a cada 5~ms, com espera baseada em \codigo{vTaskDelayUntil}. Sinais determinísticos simulam giroscópio e acelerômetro. Para roll e pitch, o filtro complementar combina a integração do giroscópio com a estimativa de aceleração:
\begin{equation}
\widehat{\theta}_k=\alpha(\widehat{\theta}_{k-1}+\omega_k\Delta t)+(1-\alpha)\theta_{a,k},
\qquad \alpha=0{,}98,\quad\Delta t=0{,}005\,\mathrm{s}.
\end{equation}
Yaw é integrado e limitado ao intervalo angular adotado. O estado contém número de sequência e timestamp de amostragem, sendo publicado por \codigo{xQueueOverwrite} em uma fila de uma posição.

A \codigo{CTRL_ATT} bloqueia nessa fila até receber um estado. O PID de cada eixo calcula a contribuição proporcional, integral e derivativa; a integral é limitada e a saída alimenta a mistura de quatro motores. Os comandos são restringidos entre 0\% e 100\%, com referência de sustentação simulada de 50\%. O modelo permite exercitar o encadeamento temporal, mas não corresponde a uma validação de estabilidade de um drone físico.

A \codigo{NAV_PLAN} processa os eventos B e C em uma fila comum. B alterna o waypoint e atualiza referências de atitude. C obtém uma fotografia do estado e imprime sequência, atitude, motores e condições de estresse/segurança. Seções críticas curtas protegem as estruturas compartilhadas.

A \codigo{FAIL_SAFE}, equivalente à \codigo{FS_TASK} do enunciado, recebe D em uma fila própria, ativa uma condição segura persistente e zera os quatro motores simulados. O CTRL consulta essa condição antes de escrever nos atuadores, impedindo a reativação após a emergência. O fim instrumentado de FS ocorre depois de uma carga artificial adicional, característica relevante para interpretar sua resposta.

As tarefas \codigo{METRICS} e \codigo{EXP_GUIDE}, de prioridade 1, emitem os resumos e orientam os toques. Em UNICORE, também executam no núcleo 0 e participam da carga global, embora não apareçam como categorias próprias nas quatro métricas principais.

\subsection{Touch, interrupções e estresse}

Foi utilizada a API \codigo{driver/touch_sens.h}, compatível com o ambiente instalado. A documentação e o exemplo oficial descrevem criação do controlador/canais, leitura e registro de callbacks~\cite{idf_touch,idf_touch_example}. O mapeamento escolhido é mostrado na Tabela~\ref{tab:touch}.

\begin{table}[!htbp]\centering\small
\caption{Entradas touch utilizadas.}\label{tab:touch}
\begin{tabularx}{\linewidth}{lllY}\toprule
Entrada & Canal & Pino & Função\\\midrule
A & T0 & GPIO4 & Perturbação e carga adicional da fusão\\
B & T7 & GPIO27 & Comando de rota/waypoint\\
C & T4 & GPIO13 & Solicitação de telemetria\\
D & T9 & GPIO32 & Emergência e condição segura\\\bottomrule
\end{tabularx}
\par\smallskip{\footnotesize Fonte: configuração da aplicação; correspondência canal/GPIO conforme documentação do ESP32~\cite{idf_touch}.}
\end{table}

Após a calibração com os contatos livres, o limite de ativação corresponde a 80\% da referência lida. A aplicação aceita um evento por contato, com debounce de 1.200~ms por canal e liberação estável de 500~ms. O callback registra o instante da detecção e encaminha o evento por operação apropriada à ISR. Assim, o trabalho extenso é realizado nas tarefas.

A carga de cálculo é reproduzida por \codigo{simulate_fixed_work}, com laço fixo e acumulador volátil. A Tabela~\ref{tab:carga} registra os parâmetros consultados. A carga tem quantidade de operações controlada; a duração depende da frequência e da interferência.

\begin{table}[!htbp]\centering\small
\caption{Quantidade de trabalho artificial na versão consultada.}\label{tab:carga}
\begin{tabularx}{\linewidth}{lrY}\toprule
Trecho & Iterações & Condição\\\midrule
FUS & 3.000 & Cada job\\
CTRL & 2.000 & Cada estado consumido\\
NAV & 24.000 & Caminho de atualização de rota\\
FS & 6.000 & Após zerar motores\\
Estresse da FUS & 500 & A cada quatro amostras durante o estresse\\\bottomrule
\end{tabularx}
\par\smallskip{\footnotesize Fonte: constantes de \codigo{main.c}.}
\end{table}

O touch A perturba os sinais e mantém o estresse por 2.000 amostras, nominalmente 10~s a 200~Hz. Esse intervalo é contado por amostras, não por um temporizador independente. Os custos sugeridos pelo enunciado são referências a medir: não foram pressupostos como WCET real, nem se ajustou artificialmente o resultado medido a esses valores.

\subsection{Períodos, deadlines e criticidade}

\begin{table}[!htbp]\centering\small
\caption{Modelo temporal adotado.}\label{tab:modelo}
\begin{tabularx}{\linewidth}{lP{3.1cm}P{3.5cm}Y}\toprule
Tarefa & Ativação & Deadline & Classe\\\midrule
FUS & Periódica; T=5 ms & 5 ms após liberação esperada & Hard\\
CTRL & Fila da FUS; taxa nominal até 200 Hz & 5 ms desde a amostra & Hard\\
NAV & Eventos B/C & 20 ms desde detecção na ISR & Soft\\
FS & Evento D via ISR & 10 ms desde detecção na ISR & Hard\\
Auxiliares & Relatórios e guia & Sem deadline funcional de voo & Sem requisito crítico\\\bottomrule
\end{tabularx}
\par\smallskip{\footnotesize Fonte: elaboração do autor conforme os requisitos e as referências temporais do firmware.}
\end{table}

FUS e CTRL foram classificadas como hard porque mantêm a cadeia de atualização dos atuadores; atrasos comprometem o ciclo temporal pretendido. FS representa a segurança e também é hard. NAV foi classificada como soft, pois atrasar uma referência de alto nível ou uma consulta não equivale à interrupção imediata do controle local. A classificação descreve a consequência de perder um prazo; não é alterada para soft quando o ensaio apresenta falhas.

O timestamp usado por CTRL é capturado no início da FUS. Dessa forma, o deadline do controle abrange fusão, espera e processamento do controlador. As tarefas auxiliares não têm requisito funcional crítico, embora possam afetar a execução das demais.

\subsection{Políticas e configuração do escalonador}

\begin{table}[!htbp]\centering\small
\caption{Prioridades implementadas; maior número indica maior prioridade.}\label{tab:prioridades}
\begin{tabularx}{\linewidth}{lrrrrY}\toprule
Política & FUS & CTRL & FS & NAV & Critério\\\midrule
CUSTOM & 4 & 3 & 5 & 2 & Segurança, fusão, controle e navegação\\
DM & 5 & 5 & 4 & 3 & Deadlines crescentes: 5, 5, 10 e 20 ms\\
RM adaptado & 5 & 5 & 4 & 3 & Cadeia nominal de 200 Hz acima dos eventos; desempate por segurança\\\bottomrule
\end{tabularx}
\par\smallskip{\footnotesize Fonte: seleção por \codigo{PRIORITY_POLICY} em \codigo{main.c}.}
\end{table}

CUSTOM privilegia a emergência, aceitando que ela interfira no controle. DM protege primeiro os menores deadlines. FUS e CTRL empatam em DM/RM, tornando o time slicing potencialmente relevante quando ambas estão prontas. As prioridades implementadas de RM e DM coincidem; por isso, diferenças entre suas execuções não representam uma ordenação diferente.

Para RM, o debounce por pad limita eventos individuais, mas NAV reúne B e C, que podem chegar próximos. Sua taxa agregada não possui intervalo mínimo de 1.200~ms. Trata-se de uma adaptação experimental de RM, sem uma demonstração teórica para chegadas agregadas e tarefas encadeadas.

O menuconfig selecionou UNICORE e frequência. Os valores de \codigo{configUSE_PREEMPTION} e \codigo{configUSE_TIME_SLICING} foram alterados na configuração do FreeRTOS para as coletas. O escalonamento por prioridades, o rodízio entre prioridades iguais e os pontos de bloqueio/yield são mecanismos distintos~\cite{barry_freertos}. No caminho touch do SDK consultado, o callback pode solicitar \codigo{portYIELD_FROM_ISR}~\cite{idf_touch_driver}; portanto, os ensaios com preempção zero são descritos como configuração cooperativa com yields explícitos, inclusive na saída da ISR.

\FloatBarrier
